\chapter{Discussion}
\label{ch:discussion}

This chapter interprets the findings, relates them to the literature, discusses what they mean for practice, and sets out the limitations that the reader should keep in mind.

\section{What the thesis finds}
\label{sec:disc-findings}

\paragraph{Archetypes and transitions.} The first supporting question asked which archetypes a robust mixed-data fuzzy clustering identifies in the Primeira Liga and how players move between them. Four archetypes of outfield players are identified, defenders, central midfielders, wide attackers and strikers, together with a small noise cluster of high-volume attackers whose output resembles none of them. The archetypes are recognisable from their medoids and statistics, and they persist over eight seasons: the same types of players are the medoids, and 75\% of players keep their archetype from one season to the next. Persistence differs greatly by archetype. Defenders are almost fixed (88\% remain), while wide attackers (54\%), and to a lesser extent strikers (65\%) and the noise cluster (25\%), are volatile, and movements occur mostly between adjacent roles, from wide attackers to central midfielders or strikers, and never between defenders and strikers. Part of the volatility of attacking players may reflect genuine changes of role or form, and part of it noise in the memberships, which the data cannot separate.

\paragraph{Explanatory power for changes in value.} The main question was whether membership transitions have incremental explanatory power for changes in market value beyond individual performance metrics and basic characteristics. In the main configuration the answer is a qualified yes. The membership changes are jointly significant against the full baseline (bootstrap $p=0.015$) and borderline against the core baseline ($p=0.079$); they raise the adjusted $R^2$ by 0.3 to 0.6 percentage points and the out-of-sample $R^2$ by 0.1 to 0.3 points; the association is present with player fixed effects and under most variations of the sample and the outcome; and a linear summary of the same statistics, the principal components, shows none. The association is concentrated in one movement, from the defender profile to the wide-attacker profile, which is associated with a smaller increase in value, by about 2 to 3 per cent for a shift of 0.10 in membership.

The qualification is considerable. First, the association depends on the clustering. It is absent for near-crisp memberships and for five archetypes, strongest for moderately soft memberships, and weaker when fewer players are assigned to the noise cluster. Second, it depends on the random draw: in 30 repeated clusterings of the main configuration, the full-baseline test is significant in 60\% of the draws and the core-baseline test in 37\%, with $p$-values that range over several orders of magnitude, and the run reported in the Results is more favourable than a typical draw. Third, the improvement in out-of-sample prediction that the memberships bring is found against the full baseline, which predicts worse than the core baseline because of its twenty-four additional statistics, but in a typical draw not against the core baseline itself. Fourth, no result survives a correction for the number of comparisons. The most accurate summary is that there is weak and specification-dependent evidence of a small association between the movement from a defender to a wide-attacker profile and changes in market value, and that the thesis does not provide evidence that membership transitions improve the prediction of changes in value.

\section{Interpreting the dependence on the clustering}
\label{sec:disc-fuzziness}

The most informative finding of the sensitivity analysis is how the association depends on the fuzziness exponent. Three observations suggest a reading, which remains a conjecture because it was not tested directly. First, with $m=1.2$ the memberships are close to crisp, since 76\% of players have a largest membership of at least 0.7, so that a transition is mostly either a jump between archetypes or no change. Such a variable discards the graded information on where a player lies between archetypes and, since a small change in the data can move a player over the boundary, is also a noisy one, which attenuates its coefficient; no association is found. Second, at $m=1.4$ and $1.5$ the memberships are graded (37 to 49\% have a largest membership of at least 0.7) while the archetypes are still well defined, and changes in them carry more of the continuous movement of a player in the space of statistics; the association is strongest here. Third, at $m=2.0$ the memberships are almost uniform (4\% at least 0.7), they hardly describe archetypes any more, and the association weakens. That the principal component scores of the same statistics show no association supports the view that what the memberships capture is not a linear change in the underlying statistics, but the relative position of a player with respect to the archetypes.

The variation across random draws, in turn, shows that this reading cannot be pressed too far: a different draw at the same $m$ changes the $p$-value by orders of magnitude. The memberships are estimated with substantial noise, and the part of the association that survives across settings is better described by its median over draws than by any single run.

The substantive content of the association is limited. In every configuration for which coefficients are shown, all coefficients are negative relative to the defender profile, and the one that stands out is that of the wide attackers: $-0.22$ against the core baseline and $-0.31$ against the full baseline at $m=1.3$, rising to $-0.31$ and $-0.41$ at $m=1.5$. A shift of 0.10 in membership is associated with a change in value of about $-2$ to $-4$ per cent, to be compared with a standard deviation of the change in log value of 0.63. Several explanations are possible, and the data in this thesis cannot discriminate between them. Valuations may depend on the role in ways that a baseline with a position indicator and a few statistics does not capture, for instance through different age profiles of value for attackers and defenders; a deviation from a stable, defensive profile towards a wide attacking one may signal a less settled role; or a large attacking output in one season may regress towards the mean in the next, which would link movements towards the wide-attacker profile to subsequent lower valuations. None of these has been examined here.

\section{Relation to the literature}
\label{sec:disc-literature}

The thesis extends the work of \citeA{durso2022robust} in the directions that they themselves suggest, by adding the temporal dimension and a financial variable. It also contains a methodological result that is relevant for anyone applying their model. The estimation of the attribute weights, which is the distinguishing feature of the model, breaks down when the groups of variables differ in dimension and in the number of distinct values, as they do in this and in most football data sets, because it favours the most compact group. The weights then had to be fixed in advance, in line with \citeA{akhanli2023clustering}. The alignment of the archetypes across seasons also needed care: matching them by their medoids, which are single players, gave inconsistent labels in most clustering runs, whereas matching them by their profiles did not. In contrast to \citeA{pacifico2025fuzzy}, who works with a panel of 27 players, the longitudinal analysis here covers all outfield players of a league over eight seasons, with a link to market values, and the finding that transitions are only weakly related to market values is a qualification of the suggestion that monitoring cluster transitions would be informative for practitioners. The baseline results are in line with the valuation literature \cite{franceschi2024determinants}: value growth falls with age, increases with playing time and with recent form of the team, and increases with goals per 90 minutes, while position and a change of club within the league have no significant association in the baseline. The thesis responds to the call of \citeA{franceschi2024determinants} to study leagues outside the big five, with the Primeira Liga as the example they name, and it confirms the concern of \citeA{muller2017beyond} that crowd valuations are noisy for smaller leagues, since in about 14\% of the pairs the value does not change from one season to the next.

\section{Implications for practice}
\label{sec:disc-practice}

For a club or an agent, the findings suggest caution. The archetypes provide a compact and interpretable description of the profile of a player and of how it changes, which may be useful for scouting and for monitoring the development of players, and the memberships of players in them show that profiles are fuzzy and change. As a signal for valuation or for identifying undervalued players, however, membership transitions cannot be recommended on this evidence: where an association is found, it is small, it depends on the tuning of the clustering and on the random draw, and it does not reliably improve prediction over a baseline of age, playing time, recent form and goals and assists, which is a more important determinant of changes in value than any of the transition variables. A practitioner who wants to use the approach should treat the memberships as noisy estimates, repeat the clustering with different seeds, report results for several settings, and not interpret a single membership change as a firm signal.

\section{Limitations}
\label{sec:disc-limitations}

\paragraph{Measurement of the profile.} FBref does not publish expected goals, expected assists or passing and ball-progression statistics for the Primeira Liga over this period, so the archetypes describe attacking output, defensive actions and positions but not the way in which players move the ball. Richer data would give richer archetypes and could change the transitions. Positions are coarse, with three labels and five common combinations, and the weights of the attribute groups were fixed instead of estimated, which is a departure from the model of \citeA{durso2022robust} that was made for the reasons given in Chapter~\ref{ch:methodology}. The noise distance follows \citeA{dave1991noise} and not the printed formula of \citeA{durso2022robust}, and the shrinkage of rates is an additional step. Each departure was necessary for a well-behaved solution on these data, but each may matter for the results.

\paragraph{Estimated regressors and random draws.} Memberships are noisy estimates; start stability is about 0.5, the bootstrap standard errors are on average 19\% larger than the conventional ones, and the bootstrap estimates are shifted towards zero relative to the original coefficients. The bootstrap corrects the standard errors but not the attenuation, so effects may be understated, and there is no measure of by how much. Because the clustering is stochastic, a single run is one draw; the 30-draw analysis is done for the main configuration at five fuzziness values, not for every specification, and the bootstrap, which takes the original run as given, does not include the dependence of the original estimate on its draw.

\paragraph{Specification search.} The conclusion depends on tuning parameters, and the thesis explored many specifications after the main configuration had been fixed. The main configuration is reported as the pre-specified one, and all specifications are reported, but the results at other settings, in particular $m=1.4$ and $1.5$, have the status of exploratory findings. The proposal also envisaged selecting $m$ with the Xie--Beni index, which was not done; the value that it would have chosen within the main configuration, $m=1.4$, gives somewhat stronger results. The pipeline was also revised during the work: the alignment of the archetypes was corrected after a check showed inconsistent labels, and a data error in the merging of the FBref tables was repaired, each time because of a diagnostic and not an outcome, but the results of the earlier versions differed, which illustrates the sensitivity discussed above. No correction for the number of specifications leaves any result significant.

\paragraph{Data and sample.} The valuation stage uses 1,361 pairs of 697 players, and the power to detect small effects is limited. Only players who stay in the league with sufficient playing time are observed in two consecutive seasons, so the findings do not apply to players who leave, who may be the ones whose value changes most. Market values are crowd estimates that move in steps, thinner for smaller leagues, and about 14\% of the pairs show no change. The baseline omits contract length and international appearances, which are strong determinants of value in the literature \cite{franceschi2024determinants,poli2024statistical}, and it does not model the transfer fee. The link between FBref and Transfermarkt is made by name and birth year, with 96\% of the player-seasons linked, a few wrong links removed by hand, and two players of the sample with the same name merged by FBref.

\paragraph{Design.} Independent clustering of the seasons ignores temporal dependence in the profiles; the number of ambiguous alignments is small at four archetypes (2 of 28 matches), and the alignment by profiles is consistent in all 120 runs, but the sensitivity of the results to the alignment method has not been examined beyond the comparison with the medoid-based version. The design is observational, and all findings are associations.

\section{Directions for future research}
\label{sec:disc-future}

Four directions follow. First, richer data, either from a provider with ball-progression and passing statistics or from event data, would allow archetypes that include how players move the ball. Second, the larger leagues provide many more transitions, and thus more power to detect small associations; a replication there, with the tuning decided in advance and the clustering repeated with many seeds, would show whether the dependence on fuzziness and on the draw is a feature of the Primeira Liga. Third, the noise in the memberships could be addressed directly, for example by methods that correct for errors in the regressors, by averaging the memberships over many random draws, or by models that estimate the profile and its change jointly instead of in two stages. Fourth, the association could be examined on longer horizons than one season, with contract length and international status among the controls, and for transfer fees in addition to market values.
